\section{Regras de inferência derivadas}\label{sec:maisSobreDemonstracoesFormais}
\index{teoria da prova}
\index{regra de inferência}
\index{teorema!da dedução}

Na seção anterior, definimos demonstrações formais em teorias de primeira ordem.
A definição de demonstração formal é bastante rígida: é uma sequência de fórmulas, cada uma das quais é um axioma lógico, um axioma da teoria, ou é obtida a partir de fórmulas anteriores por uma das regras de inferência explicitamente permitidas.

Na prática, demonstrações matemáticas são escritas de forma muito mais flexível.
Escreve-se, a todo momento, expressões como ``suponha que $\varphi$'', ``fixe $x$ arbitrário'', ``tome $x$ tal que $\varphi(x)$'', ``dividimos em casos'' ou ``basta provar que $\psi$'', dentre muitas outras formas de raciocínio que não estão explicitamente previstas na definição de demonstração formal.
O objetivo desta seção é justificar, a partir desta noção formal, várias das principais formas de raciocínio usadas na prática.

As regras que aparecerão nesta seção não são novas regras primitivas acrescentadas ao sistema.
São regras derivadas: cada uma delas é um metateorema que, dadas demonstrações formais das premissas indicadas, fornece um procedimento finitário para construir uma demonstração formal da conclusão.

\subsection{Notação para regras de inferência}

Uma regra de inferência será descrita pela notação a seguir.
\[
    \frac{\text{Premissas}}{\text{Conclusão}}
\]
Essa notação significa que, sempre que dispomos dos dados indicados nas premissas, há um procedimento finito que constrói a demonstração formal indicada na conclusão.

Por exemplo, o Metalema~\ref{mlem:modusPonens} (Modus Ponens) pode ser escrito como a regra
\[
    \frac{T\vdash \varphi \quad T\vdash \varphi\rightarrow\psi}
         {T\vdash \psi}.
\]
Tal notação nos diz que a partir de demonstrações formais de $\varphi$ e de $\varphi\rightarrow\psi$ em $T$, é possível construir uma demonstração formal de $\psi$ em $T$.

Utilizaremos a convenção de que, se $T=(L,\Sigma)$ é uma teoria de primeira ordem e $\varphi$ é uma $L$-sentença, então $T\cup \{\varphi\}$ é a teoria $(L,\Sigma\cup\{\varphi\})$.

\subsection{O Teorema da Dedução}

O Teorema da Dedução formaliza o uso de hipóteses temporárias.
Ele justifica o procedimento usual de, para provar uma implicação $\varphi\rightarrow\psi$ em que $\varphi$ é uma $L$-sentença, assumir temporariamente $\varphi$ e, sob essa hipótese, provar $\psi$.

\begin{metatheorem}[Teorema da Dedução]\label{mteo:teoremaDaDeducaoFO}\label{mthm:teoremaDaDeducao}
    Seja $T=(L,\Sigma)$ uma teoria de primeira ordem, $\varphi$ uma $L$-sentença e seja $\psi$ uma $L$-fórmula.

    Então $T\cup\{\varphi\}\vdash \psi$ se, e somente se, $T\vdash \varphi\rightarrow\psi$.
    Ou seja:
    \[
        \frac{T\cup\{\varphi\}\vdash \psi}
             {T\vdash \varphi\rightarrow\psi} \quad \text{e}\quad \frac{T\vdash \varphi\rightarrow\psi}
             {T\cup\{\varphi\}\vdash \psi}.
    \]
\end{metatheorem}

\begin{proof}
    Suponha primeiramente que $T\vdash \varphi\rightarrow\psi$.
    Da monotonicidade da dedução, segue que $T\cup\{\varphi\}\vdash \varphi\rightarrow\psi$.
    Note que $T\cup\{\varphi\}\vdash \varphi$.
    Dessa forma, $T\cup\{\varphi\}\vdash \psi$ por modus ponens.
    
    Reciprocamente, suponha que $T\cup\{\varphi\}\vdash \psi$.
    Fixe uma demonstração formal de $\psi$ a partir de $T\cup\{\varphi\}$, digamos $\psi_0,\dots,\psi_n$, em que $\psi_n$ é $\psi$.
    Provaremos, por indução em $i\leq n$, que
    \[
        T\vdash \varphi\rightarrow\psi_i.
    \]

    Suponha que a tese já tenha sido provada para todo $j<i$.
    Veremos que $T\vdash \varphi\rightarrow\psi_i$.

    \textbf{Caso 1}: $\psi_i$ é um axioma lógico ou um axioma de $T$.
    Nesse caso, $T\vdash \psi_i$.
    Ademais, $\psi_i\rightarrow(\varphi\rightarrow\psi_i)$ é uma tautologia (ou uma instância de FO1), de modo que $T\vdash \psi_i\rightarrow(\varphi\rightarrow\psi_i)$.
    Por modus ponens, segue que $T\vdash \varphi\rightarrow\psi_i$.
    
    \textbf{Caso 2}: $\psi_i$ é a fórmula $\varphi$.
    Nesse caso, note que $T\vdash \varphi\rightarrow\varphi$, pois esta fórmula é uma instância de tautologia.

    \textbf{Caso 3}: $\psi_i$ foi obtida por modus ponens a partir de fórmulas anteriores.
    Então existem $j,k<i$ tais que $\psi_j$ é $\psi_k\rightarrow\psi_i$.
    Pela hipótese de indução,
    \[
        T\vdash \varphi\rightarrow\psi_k \quad \text{e}\quad T\vdash \varphi\rightarrow(\psi_k\rightarrow\psi_i).
    \]
    Por FO2 (ou porque tal fórmula é uma instância de tautologia), temos
    \[
        T\vdash (\varphi\rightarrow(\psi_k\rightarrow\psi_i))
        \rightarrow
        ((\varphi\rightarrow\psi_k)\rightarrow(\varphi\rightarrow\psi_i)).
    \]
    Duas aplicações de modus ponens concluem que $T\vdash \varphi\rightarrow\psi_i$.

    \textbf{Caso 4}: suponha que $\psi_i$ foi obtida pela regra de introdução do quantificador universal.
    Então existem fórmulas $\chi$ e $\eta$, uma variável $x$ e $j<i$ tais que $\psi_j$ é $\chi\rightarrow\eta$, $\psi_i$ é $\chi\rightarrow\forall x\,\eta$, e $x$ não ocorre livre em $\chi$.

    Pela hipótese de indução, $T\vdash \varphi\rightarrow(\chi\rightarrow\eta)$.
    Temos que $T\vdash (\varphi\rightarrow(\chi\rightarrow\eta))\rightarrow((\varphi\land\chi)\rightarrow\eta)$, pois tal fórmula é uma instância de tautologia.
    Assim, por modus ponens, segue que
    \[
        T\vdash (\varphi\land\chi)\rightarrow\eta.
    \]
    Como $\varphi$ é uma sentença e $x$ não ocorre livre em $\chi$, a variável $x$ não ocorre livre em $\varphi\land\chi$.
    Pela regra de introdução do quantificador universal,
    \[
        T\vdash (\varphi\land\chi)\rightarrow\forall x\,\eta.
    \]
    Usando a instância de tautologia $((\varphi\land\chi)\rightarrow\forall x\,\eta)
        \rightarrow
        (\varphi\rightarrow(\chi\rightarrow\forall x\,\eta))$ e modus ponens, obtemos
    \[
        T\vdash \varphi\rightarrow(\chi\rightarrow\forall x\,\eta),
    \]
    isto é, $T\vdash \varphi\rightarrow\psi_i$.

    \textbf{Caso 5}: $\psi_i$ foi obtida pela regra de introdução do quantificador existencial.
    Então existem fórmulas $\chi$ e $\eta$, uma variável $x$ e $j<i$ tais que $\psi_j$ é $\chi\rightarrow\eta$, $\psi_i$ é $\exists x\chi\rightarrow\eta$, e $x$ não ocorre livre em $\eta$.

    Pela hipótese de indução, $T\vdash \varphi\rightarrow(\chi\rightarrow\eta)$.
    Pela instância de tautologia $(\varphi\rightarrow(\chi\rightarrow\eta))\rightarrow(\chi\rightarrow(\varphi\rightarrow\eta))$ e modus ponens, obtemos
    \[
        T\vdash \chi\rightarrow(\varphi\rightarrow\eta).
    \]
    Como $\varphi$ é uma sentença e $x$ não ocorre livre em $\eta$, a variável $x$ não ocorre livre em $\varphi\rightarrow\eta$.
    Pela regra de introdução do quantificador existencial,
    \[
        T\vdash \exists x\chi\rightarrow(\varphi\rightarrow\eta).
    \]
    Usando a instância de tautologia $(\exists x\chi\rightarrow(\varphi\rightarrow\eta))
        \rightarrow
        (\varphi\rightarrow(\exists x\chi\rightarrow\eta))$ e modus ponens, concluímos
    \[
        T\vdash \varphi\rightarrow(\exists x\chi\rightarrow\eta),
    \]
    isto é, $T\vdash \varphi\rightarrow\psi_i$.

    Portanto, $T\vdash \varphi\rightarrow\psi_i$ para todo $i\leq n$.
    Como $\psi_n$ é $\psi$, segue que $T\vdash \varphi\rightarrow\psi$.
\end{proof}

Em linguagem mais informal, o Teorema da Dedução autoriza o seguinte raciocínio corriqueiro em matemática, quando $\varphi$ é uma sentença:

``Queremos mostrar que $\varphi\rightarrow \psi$.
Suponha $\varphi$.
Veremos que $\psi$.''

\subsection{Regras proposicionais derivadas}

Começaremos com regras que dependem apenas dos conectivos proposicionais.
Elas são consequências diretas dos axiomas FO1--FO13, do fato de que instâncias de tautologias são demonstráveis e de \emph{modus ponens}.

\begin{metaproposition}[Regras proposicionais para implicação]\label{mprop:regrasProposicionaisImplicacao}
    Seja $T$ uma teoria de primeira ordem e sejam $\varphi,\psi,\theta$ fórmulas da linguagem de $T$.

    Então valem as seguintes regras de inferência.
    Abaixo do enunciado de cada uma delas, escreveremos uma forma informal corriqueira de expressar o raciocínio que a regra justifica.

    \begin{enumerate}[label=(\roman*)]
        \item Modus ponens.
        \[
            \frac{T\vdash \varphi \quad T\vdash \varphi\rightarrow\psi}
                 {T\vdash \psi}.
        \]
        
        ``Como vale $\varphi$, e $\varphi$ implica $\psi$, temos $\psi$.''

        \item Modus tollens.
        \[
            \frac{T\vdash \varphi\rightarrow\psi \quad T\vdash \neg\psi}
                 {T\vdash \neg\varphi}.
        \]

        ``Como vale $\neg\psi$, e $\varphi$ implica $\psi$, temos $\neg\varphi$.''

        \item Contrapositiva.
        
        \[
            \frac{T\vdash \varphi\rightarrow\psi}
                 {T\vdash \neg\psi\rightarrow\neg\varphi} \qquad \frac{T\vdash \neg\psi\rightarrow\neg\varphi}
                 {T\vdash \varphi\rightarrow\psi}
        \]

        ``Queremos mostrar que $\varphi$ implica $\psi$.
        Procedemos pela contrapositiva mostrando que $\neg\psi$ implica $\neg\varphi$.''

        \item Introdução de $\rightarrow$.
        \[
            \frac{T\vdash \psi}
                 {T\vdash \varphi\rightarrow\psi} \qquad \frac{T\vdash \neg\varphi}
                 {T\vdash \varphi\rightarrow\psi}
        \]

        ``Como $\psi$ é verdadeiro, $\varphi$ implica $\psi$ independentemente de $\varphi$.''

        ``Como o antecedente $\varphi$ é falso, $\varphi$ implica $\psi$ trivialmente.''


        \item Reflexividade de $\rightarrow$.
        \[
            T\vdash \varphi\rightarrow\varphi.
        \]

        ``Trivialmente, $\varphi$ implica $\varphi$.''

        \item Transitividade de $\rightarrow$.
        \[
            \frac{T\vdash \varphi\rightarrow\psi \quad T\vdash \psi\rightarrow\theta}
                 {T\vdash \varphi\rightarrow\theta}.
        \]

        ``Como $\varphi$ implica $\psi$ e $\psi$ implica $\theta$, temos que $\varphi$ implica $\theta$.''
    \end{enumerate}
\end{metaproposition}

\begin{proof}
    (i) é exatamente o Metalema~\ref{mlem:modusPonens}.
    Para os demais itens, utilizamos modus ponens com as seguintes instâncias de tautologia.

    (ii) $(\varphi\rightarrow\psi)\rightarrow(\neg\psi\rightarrow\neg\varphi)$

    (iii) $(\neg\psi\rightarrow\neg\varphi)\rightarrow(\varphi\rightarrow\psi)$ e $(\varphi\rightarrow\psi)\rightarrow(\neg\psi\rightarrow\neg\varphi)$

    (iv) $\psi\rightarrow(\varphi\rightarrow\psi)$ e $\neg\varphi\rightarrow(\varphi\rightarrow\psi)$

    (v) $\varphi\rightarrow\varphi$

    (vi) $(\varphi\rightarrow\psi)\rightarrow((\psi\rightarrow\theta)\rightarrow(\varphi\rightarrow\theta))$
\end{proof}

\begin{metaproposition}[Regras proposicionais para negação]\label{mprop:regrasProposicionaisNegacao}
    Seja $T$ uma teoria de primeira ordem e sejam $\varphi$, $\psi$ fórmulas da linguagem de $T$.

    Então valem as seguintes regras de inferência.
    Abaixo do enunciado de cada uma delas, escreveremos uma forma informal corriqueira de expressar o raciocínio que a regra justifica.

    \begin{enumerate}[label=(\roman*)]
        \item Eliminação da dupla negação.
        \[
            \frac{T\vdash \neg\neg\varphi}
                 {T\vdash \varphi}.
        \]
        
        ``Como não vale $\neg \varphi$, temos $\varphi$.''

        \item Redução ao absurdo (Reductio ad absurdum). Se $\varphi$ é uma sentença, então
        \[
            \frac{T\cup\{\varphi\}\vdash \psi \quad T\cup\{\varphi\}\vdash \neg\psi}
                 {T\vdash \neg\varphi}.
        \]

        ``Suponha por absurdo que vale $\varphi$ (...) como $\psi$ e $\neg\psi$, temos um absurdo.
        Logo, $\neg\varphi$.''
    \end{enumerate}
\end{metaproposition}

\begin{proof}
    Para (i) basta observar que $\neg\neg\varphi\rightarrow\varphi$ é um axioma lógico (ou uma instância de tautologia) e utilizar modus ponens.

    Para (ii), pelo Teorema da Dedução, $T\vdash \varphi\rightarrow\psi$ e $T\vdash \varphi\rightarrow\neg\psi$.
    Por FO3,
    \[
        T\vdash
        (\varphi\rightarrow\psi)
        \rightarrow
        ((\varphi\rightarrow\neg\psi)\rightarrow\neg\varphi).
    \]
    Aplicando modus ponens duas vezes, obtemos $T\vdash \neg\varphi$.
\end{proof}

\begin{metaproposition}[Regras proposicionais para conjunção]\label{mprop:regrasProposicionaisConjuncao}
    Seja $T$ uma teoria de primeira ordem e sejam $\varphi$, $\psi$ e $\theta$ fórmulas da linguagem de $T$.

    Então valem as seguintes regras de inferência.
    Abaixo do enunciado de cada uma delas, escreveremos uma forma informal corriqueira de expressar o raciocínio que a regra justifica.
    \begin{enumerate}[label=(\roman*)]

        \item Introdução de $\land$.
        \[
            \frac{T\vdash \varphi \quad T\vdash \psi}
                 {T\vdash \varphi\land\psi}.
        \]

        ``Como temos $\varphi$ e $\psi$, temos $\varphi\land\psi$.''

        \item Eliminação de $\land$.
        \[
            \frac{T\vdash \varphi\land\psi}
                 {T\vdash \varphi}
            \qquad
            \frac{T\vdash \varphi\land\psi}
                 {T\vdash \psi}.
        \]

        ``Temos $\varphi\land\psi$, logo, $\varphi$ e $\psi$.''

        \item Princípio da não contradição.
        \[
            T\vdash \neg(\varphi\land\neg\varphi).
        \]

        ``$\varphi$ e $\neg\varphi$ não podem ambos ser verdadeiros.''

        \item Princípio da explosão.
        \[
            T\vdash (\varphi\land\neg\varphi)\rightarrow\psi.
        \]

        ``Como temos $\varphi$ e $\neg\varphi$, temos uma contradição.
        De uma contradição, qualquer fórmula $\psi$ pode ser provada.
        Logo, segue $\psi$.''

        \item Comutatividade de $\land$.
        \[
            \frac{T\vdash \varphi\land\psi}
                 {T\vdash \psi\land\varphi}.
        \]
        ``Temos $\varphi\land\psi$, equivalentemente, $\psi\land\varphi$.''

        \item Associatividade de $\land$.
        \[
            \frac{T\vdash (\varphi\land\psi)\land\theta}
                 {T\vdash \varphi\land(\psi\land\theta)}
            \qquad
            \frac{T\vdash \varphi\land(\psi\land\theta)}
                 {T\vdash (\varphi\land\psi)\land\theta}.
        \]

        ``Temos $(\varphi\land\psi)\land\theta$, equivalentemente, $\varphi\land(\psi\land\theta)$.
        Omitindo os parênteses, vamos escrever $\varphi\land\psi\land\theta$.''
    \end{enumerate}
\end{metaproposition}

\begin{proof}
    Tais regras de inferência decorrem de modus ponens e das tautologias $(\varphi\land \psi)\rightarrow\varphi$, $(\varphi\land \psi)\rightarrow\psi$, $\varphi\rightarrow(\psi\rightarrow(\varphi\land\psi))$, $\neg(\varphi\land\neg\varphi)$, $(\varphi\land\neg\varphi)\rightarrow\psi$, $(\varphi\land\psi)\rightarrow(\psi\land\varphi)$, $((\varphi\land\psi)\land\theta)\rightarrow(\varphi\land(\psi\land\theta))$ e $(\varphi\land(\psi\land\theta))\rightarrow((\varphi\land\psi)\land\theta)$.
\end{proof}

\begin{metaproposition}[Regras proposicionais para disjunção]\label{mprop:regrasProposicionaisDisjuncao}
    Seja $T$ uma teoria de primeira ordem e sejam $\varphi$, $\psi$ e $\theta$ fórmulas da linguagem de $T$.

    Então valem as seguintes regras de inferência.
    Abaixo do enunciado de cada uma delas, escreveremos uma forma informal corriqueira de expressar o raciocínio que a regra justifica.
    \begin{enumerate}[label=(\roman*)]

        \item Introdução de $\lor$.
        \[
            \frac{T\vdash \varphi}
                 {T\vdash \varphi\lor\psi}
            \qquad
            \frac{T\vdash \psi}
                 {T\vdash \varphi\lor\psi}.
        \]

        ``De $\varphi$, segue que $\varphi\lor\psi$.''

        ``Como temos $\psi$, temos $\varphi\lor\psi$.''

        \item Introdução de casos. Se $\varphi$ e $\psi$ são sentenças, então
        \[
            \frac{T\vdash \varphi\lor\psi \quad T\cup\{\varphi\}\vdash \theta \quad T\cup\{\psi\}\vdash \theta}
                 {T\vdash \theta}.
        \]

        ``Temos $\varphi\lor\psi$. Suponha que $\varphi$ seja verdadeira. Então, $\theta$ é verdadeira. Suponha que $\psi$ seja verdadeira. Então, $\theta$ é verdadeira. Logo, $\theta$ é verdadeira.''

        \item Princípio do terceiro excluído.
        \[
            T\vdash \varphi\lor\neg\varphi.
        \]

        ``Pelo princípio do terceiro excluído, ou $\varphi$ é verdadeira, ou $\neg\varphi$ é verdadeira.''

        \item Disjunção de casos. Se $\varphi$ é uma sentença, então
        \[
            \frac{T\cup\{\varphi\}\vdash \psi \quad T\cup\{\neg\varphi\}\vdash \psi}
                 {T\vdash \psi}.
        \]

        ``Queremos mostrar $\psi$.
        Dividimos em casos: suponha que $\varphi$ seja verdadeira (...) então $\psi$ é verdadeira; suponha que $\varphi$ seja falsa (...) então $\psi$ é verdadeira. Logo, $\psi$ é verdadeira.''

        \item Prova de disjunção. Se $\varphi$ é uma sentença, então
        \[
            \frac{T\cup\{\neg\varphi\}\vdash \psi}
                 {T\vdash \varphi\lor\psi}.
        \]

        ``Queremos mostrar $\varphi\lor\psi$.
        É suficiente assumir $\neg\varphi$ e mostrar que, nesse caso, $\psi$ é verdadeira.''

        \item Comutatividade de $\lor$.
        \[
            \frac{T\vdash \varphi\lor\psi}
                 {T\vdash \psi\lor\varphi}.
        \]

        ``Temos $\varphi\lor\psi$, equivalentemente, $\psi\lor\varphi$.''

        \item Associatividade de $\lor$.
        
        \[
            \frac{T\vdash (\varphi\lor\psi)\lor\theta}
                 {T\vdash \varphi\lor(\psi\lor\theta)}
            \qquad
            \frac{T\vdash \varphi\lor(\psi\lor\theta)}
                 {T\vdash (\varphi\lor\psi)\lor\theta}.
        \]

        ``Temos $(\varphi\lor\psi)\lor\theta$, equivalentemente, $\varphi\lor(\psi\lor\theta)$. Omitindo os parênteses, vamos escrever $\varphi\lor\psi\lor\theta$.''

        \item Distributividade de $\land$ sobre $\lor$.
        
        \[
            \frac{T\vdash \varphi\land(\psi\lor\theta)}
                 {T\vdash (\varphi\land\psi)\lor(\varphi\land\theta)}
            \qquad
            \frac{T\vdash (\varphi\land\psi)\lor(\varphi\land\theta)}
                 {T\vdash \varphi\land(\psi\lor\theta)}.
        \]

        ``Temos $\varphi\land(\psi\lor\theta)$, equivalentemente, $(\varphi\land\psi)\lor(\varphi\land\theta)$.''

        \item Distributividade de $\lor$ sobre $\land$.
        
        \[
            \frac{T\vdash \varphi\lor(\psi\land\theta)}
                 {T\vdash (\varphi\lor\psi)\land(\varphi\lor\theta)}
            \qquad
            \frac{T\vdash (\varphi\lor\psi)\land(\varphi\lor\theta)}
                 {T\vdash \varphi\lor(\psi\land\theta)}.
        \]

        ``Temos $\varphi\lor(\psi\land\theta)$, equivalentemente, $(\varphi\lor\psi)\land(\varphi\lor\theta)$.''

        \item Leis de De Morgan para $\land$.
        
        \[
            \frac{T\vdash \neg(\varphi\land\psi)}
                 {T\vdash \neg\varphi\lor\neg\psi}
            \qquad
            \frac{T\vdash \neg\varphi\lor\neg\psi}
                 {T\vdash \neg(\varphi\land\psi)}
        \]

        ``Temos $\neg(\varphi\land\psi)$.
        Por De Morgan, $\neg\varphi\lor\neg\psi$.''

        ``Temos $\neg\varphi\lor\neg\psi$.
        Equivalentemente, $\neg(\varphi\land\psi)$.''

        \item Leis de De Morgan para $\lor$.
        
        \[
            \frac{T\vdash \neg(\varphi\lor\psi)}
                 {T\vdash \neg\varphi\land\neg\psi}
            \qquad
            \frac{T\vdash \neg\varphi\land\neg\psi}
                 {T\vdash \neg(\varphi\lor\psi)}
        \]

        ``Temos $\neg(\varphi\lor\psi)$.
        Por De Morgan, $\neg\varphi\land\neg\psi$.''
        
        ``Temos $\neg\varphi\land\neg\psi$.
        Equivalentemente, $\neg(\varphi\lor\psi)$.''

        \item Idempotência de $\lor$.
        
        \[
            \frac{T\vdash \varphi\lor \varphi}{
                 T\vdash \varphi}
        \]

        ``Temos $\varphi\lor\varphi$.
        Em qualquer caso, $\varphi$.''
    \end{enumerate}
\end{metaproposition}

\begin{proof}
    Tais regras de inferência decorrem de modus ponens, do Teorema da Dedução e das seguintes instâncias de tautologia.

    (i) $\varphi\rightarrow(\varphi\lor\psi)$ e $\psi\rightarrow(\varphi\lor\psi)$.

    (ii) $(\varphi\lor\psi)\rightarrow((\varphi\rightarrow\theta)\rightarrow((\psi\rightarrow\theta)\rightarrow\theta))$.

    (iii) $\varphi\lor\neg\varphi$.

    (iv) $(\varphi\rightarrow\psi)\rightarrow((\neg\varphi\rightarrow\psi)\rightarrow((\varphi\lor\neg\varphi)\rightarrow\psi))$.

    (v) $(\neg\varphi\rightarrow\psi)\rightarrow(\varphi\lor\psi)$.

    (vi) $(\varphi\lor\psi)\rightarrow(\psi\lor\varphi)$.

    (vii) $((\varphi\lor\psi)\lor\theta)\rightarrow(\varphi\lor(\psi\lor\theta))$ e $(\varphi\lor(\psi\lor\theta))\rightarrow((\varphi\lor\psi)\lor\theta)$.

    (viii) $\varphi\land(\psi\lor\theta)\rightarrow((\varphi\land\psi)\lor(\varphi\land\theta))$ e $((\varphi\land\psi)\lor(\varphi\land\theta))\rightarrow(\varphi\land(\psi\lor\theta))$.

    (ix) $\varphi\lor(\psi\land\theta)\rightarrow((\varphi\lor\psi)\land(\varphi\lor\theta))$ e $((\varphi\lor\psi)\land(\varphi\lor\theta))\rightarrow(\varphi\lor(\psi\land\theta))$.

    (x) $\neg(\varphi\land\psi)\rightarrow(\neg\varphi\lor\neg\psi)$ e $(\neg\varphi\lor\neg\psi)\rightarrow\neg(\varphi\land\psi)$.

    (xi) $\neg(\varphi\lor\psi)\rightarrow(\neg\varphi\land\neg\psi)$ e $(\neg\varphi\land\neg\psi)\rightarrow\neg(\varphi\lor\psi)$.

    (xii) $(\varphi\lor\varphi)\rightarrow\varphi$.
\end{proof}

\begin{metaproposition}[Regras proposicionais para bicondicional]\label{mprop:regrasProposicionaisBicondicional}
    Seja $T$ uma teoria de primeira ordem e sejam $\varphi$ e $\psi$ fórmulas da linguagem de $T$.

    Então valem as seguintes regras de inferência.
    Abaixo do enunciado de cada uma delas, escreveremos uma forma informal corriqueira de expressar o raciocínio que a regra justifica.

    \begin{enumerate}[label=(\roman*)]
        \item Introdução de $\leftrightarrow$.
        \[
            \frac{T\vdash \varphi\rightarrow\psi \quad T\vdash \psi\rightarrow\varphi}
                 {T\vdash \varphi\leftrightarrow\psi}.
        \]

        ``Como $\varphi$ implica $\psi$ e $\psi$ implica $\varphi$, temos que $\varphi$ é equivalente a $\psi$.''

        \item Eliminação de $\leftrightarrow$.
        \[
            \frac{T\vdash \varphi\leftrightarrow\psi}
                 {T\vdash \varphi\rightarrow\psi}
            \qquad
            \frac{T\vdash \varphi\leftrightarrow\psi}
                 {T\vdash \psi\rightarrow\varphi}.
        \]

        ``Temos que $\varphi$ é equivalente a $\psi$, logo, $\varphi$ implica $\psi$ e $\psi$ implica $\varphi$.''

        \item Comutatividade de $\leftrightarrow$.
        \[
            \frac{T\vdash \varphi\leftrightarrow\psi}
                 {T\vdash \psi\leftrightarrow\varphi}.
        \]

        ``Temos que $\varphi$ é equivalente a $\psi$, equivalentemente, $\psi$ é equivalente a $\varphi$.''

        \item Associatividade de $\leftrightarrow$.
        \[
            \frac{T\vdash (\varphi\leftrightarrow\psi)\leftrightarrow\theta}
                 {T\vdash \varphi\leftrightarrow(\psi\leftrightarrow\theta)}
            \qquad
            \frac{T\vdash \varphi\leftrightarrow(\psi\leftrightarrow\theta)}
                 {T\vdash (\varphi\leftrightarrow\psi)\leftrightarrow\theta}.
        \]

        ``Temos $(\varphi\leftrightarrow\psi)\leftrightarrow\theta$, equivalentemente, $\varphi\leftrightarrow(\psi\leftrightarrow\theta)$.''
        
    \end{enumerate}
\end{metaproposition}
\begin{proof}
    Tais regras de inferência decorrem de modus ponens e das seguintes instâncias de tautologia.

    (i) $(\varphi\rightarrow\psi)\rightarrow((\psi\rightarrow\varphi)\rightarrow(\varphi\leftrightarrow\psi))$.

    (ii) $(\varphi\leftrightarrow\psi)\rightarrow(\varphi\rightarrow\psi)$ e $(\varphi\leftrightarrow\psi)\rightarrow(\psi\rightarrow\varphi)$.

    (iii) $(\varphi\leftrightarrow\psi)\rightarrow(\psi\leftrightarrow\varphi)$.

    (iv) $((\varphi\leftrightarrow\psi)\leftrightarrow\theta)\rightarrow(\varphi\leftrightarrow(\psi\leftrightarrow\theta))$ e $(\varphi\leftrightarrow(\psi\leftrightarrow\theta)) \rightarrow ((\varphi\leftrightarrow \psi)\leftrightarrow \theta)$.
\end{proof}

\subsection{Regras para quantificadores}

Passaremos agora às regras envolvendo quantificadores.
Aqui aparecem as condições mais delicadas da seção, pois é preciso controlar cuidadosamente variáveis livres e substituições válidas.

Comecemos com duas regras diretamente associadas aos axiomas FO14 e FO15.

\begin{metaproposition}[Instanciação universal e generalização existencial]\label{mprop:instanciacaoUniversalGeneralizacaoExistencial}
    Seja $T$ uma teoria de primeira ordem de léxico $L$.
    Seja $\varphi$ uma $L$-fórmula, seja $x$ uma variável e seja $t$ um $L$-termo tal que $x$ é substituível por $t$ em $\varphi$.

    Então:

    \begin{enumerate}[label=(\roman*)]
        \item Instanciação universal.
        
        \[
        \frac{T\vdash \forall x\varphi}
             {T\vdash \varphi[x/t]}
        \]
        ``Como $\varphi$ vale para todo $x$, então vale para este termo $t$.''
        \item Generalização existencial.
        \[
            \frac{T\vdash \varphi[x/t]}
             {T\vdash \exists x\varphi}
        \]
        ``Como $\varphi$ vale para o termo $t$, então existe um $x$ tal que $\varphi$.''
    \end{enumerate}
\end{metaproposition}

\begin{proof}
    Para (i), usamos FO14:
    \[
        \forall x\varphi\rightarrow\varphi[x/t].
    \]
    Como $T\vdash \forall x\varphi$, segue por \emph{modus ponens} que $T\vdash \varphi[x/t]$.

    Para (ii), usamos FO15:
    \[
        \varphi[x/t]\rightarrow\exists x\varphi.
    \]
    Como $T\vdash \varphi[x/t]$, segue por \emph{modus ponens} que $T\vdash \exists x\varphi$.
\end{proof}

\begin{metaproposition}[Negação de quantificadores]\label{mprop:negacaoQuantificadores}
    Seja $T$ uma teoria de primeira ordem, seja $\varphi$ uma fórmula da linguagem de $T$ e seja $x$ uma variável.
    Então
    \[
        T\vdash \neg\exists x\varphi \leftrightarrow \forall x\neg\varphi.
    \]
\end{metaproposition}

\begin{proof}
    Primeiro, provemos que $T\vdash \neg\exists x\varphi \rightarrow \forall x\neg\varphi$.

    Por FO15, aplicado ao termo $x$, temos $T\vdash \varphi\rightarrow\exists x\varphi$.
    Pela contrapositiva,
    \[
        T\vdash \neg\exists x\varphi\rightarrow\neg\varphi.
    \]
    Como $x$ não ocorre livre em $\neg\exists x\varphi$, pela regra de introdução do quantificador universal,
    \[
        T\vdash \neg\exists x\varphi\rightarrow\forall x\neg\varphi.
    \]

    Reciprocamente, provemos que $T\vdash \forall x\neg\varphi\rightarrow\neg\exists x\varphi$.

    Por FO14, aplicado à fórmula $\neg\varphi$ e ao termo $x$, temos
    \[
        T\vdash \forall x\neg\varphi\rightarrow\neg\varphi.
    \]
    Pela contrapositiva e pela equivalência entre $\neg\neg \varphi$ e $\varphi$, temos que $T\vdash \neg \neg\varphi\rightarrow\neg\forall x\neg\varphi$ e $T\vdash \varphi\rightarrow\neg\neg\varphi$.
    Assim, pela transitividade de $\rightarrow$,
    \[
        T\vdash \varphi\rightarrow\neg\forall x\neg\varphi.
    \]
    Como $x$ não ocorre livre em $\neg\forall x\neg\varphi$, pela regra de introdução do quantificador existencial,
    \[
        T\vdash \exists x\varphi\rightarrow\neg\forall x\neg\varphi.
    \]
    Aplicando novamente a contrapositiva, $T\vdash \neg \neg \forall x\neg\varphi\rightarrow\neg\exists x\varphi$.
    Como $T\vdash \forall x\neg\varphi\rightarrow\neg\neg\forall x\neg\varphi$, pela transitividade de $\rightarrow$,
    \[
        T\vdash \forall x\neg\varphi\rightarrow\neg\exists x\varphi.
    \]

    Pelas duas implicações, concluímos
    \[
        T\vdash \neg\exists x\varphi \leftrightarrow \forall x\neg\varphi.
    \]
\end{proof}

\begin{metaproposition}[Quantificadores e implicações]\label{mprop:quantificadoresImplicacao}
    Seja $T$ uma teoria de primeira ordem e $\varphi$, $\psi$ fórmulas da linguagem de $T$.
    Seja $x$ uma variável.
    Então:
    \begin{enumerate}[label=(\roman*)]
        \item 
        \[
        \frac{\phi\rightarrow\psi}
             {\forall x\phi\rightarrow\forall x\psi}
        \]
        ``Como para qualquer que seja $x$, $\phi$ implica $\psi$, então, se $\phi$ vale para todo $x$, então $\psi$ vale para todo $x$.''
        \item
        \[
        \frac{\phi\rightarrow\psi}
             {\exists x\phi\rightarrow\exists x\psi}
        \]
        ``Como para qualquer que seja $x$, $\phi$ implica $\psi$, então, se existe um $x$ tal que $\phi$, então existe um $x$ tal que $\psi$.''
        \item \[
        \frac{\phi\leftrightarrow\psi}
             {\forall x\phi\leftrightarrow\forall x\psi}
        \]
        ``Como para qualquer que seja $x$, $\phi$ e $\psi$ são equivalentes, então, $\phi$ vale para qualquer que seja $x$ se, e somente se, $\psi$ vale para qualquer que seja $x$.''
        \item \[
        \frac{\phi\leftrightarrow\psi}
             {\exists x\phi\leftrightarrow\exists x\psi}
        \]
        ``Como para qualquer que seja $x$, $\phi$ e $\psi$ são equivalentes, então, existe um $x$ tal que $\phi$ se, e somente se, existe um $x$ tal que $\psi$.''
    \end{enumerate}
\end{metaproposition}
\begin{proof}
    (i) Por FO14, aplicado à fórmula $\psi$ e ao termo $x$, temos $T\vdash \psi\rightarrow\forall x\psi$.
    Como $x$ não ocorre livre em $\psi\rightarrow\forall x\psi$, pela regra de introdução do quantificador universal, $T\vdash \forall x\phi\rightarrow\forall x\psi$.

    (ii) Por FO15, aplicado à fórmula $\phi$ e ao termo $x$, temos $T\vdash \exists x\phi\rightarrow\phi$.
    Como $x$ não ocorre livre em $\exists x\phi\rightarrow\phi$, pela regra de introdução do quantificador existencial, $T\vdash \exists x\phi\rightarrow\exists x\psi$.

    (iii) Se $T\vdash \phi\leftrightarrow\psi$, então $T\vdash \phi\rightarrow\psi$ e $T\vdash \psi\rightarrow\phi$.
    Por (i), $T\vdash \forall x\phi\rightarrow\forall x\psi$ e $T\vdash \forall x\psi\rightarrow\forall x\phi$.
    Assim, $T\vdash \forall x\phi\leftrightarrow\forall x\psi$.

    (iv) Se $T\vdash \phi\leftrightarrow\psi$, então $T\vdash \phi\rightarrow\psi$ e $T\vdash \psi\rightarrow\phi$.
    Por (ii), $T\vdash \exists x\phi\rightarrow\exists x\psi$ e $T\vdash \exists x\psi\rightarrow\exists x\phi$.
    Assim, $T\vdash \exists x\phi\leftrightarrow\exists x\psi$.
\end{proof}

\begin{metacorollary}[Negação de quantificadores II]\label{mprop:negacaoQuantificadores}
    Seja $T$ uma teoria de primeira ordem, seja $\varphi$ uma fórmula da linguagem de $T$ e seja $x$ uma variável.
    Então
    \[
        T\vdash \neg \forall x\varphi\leftrightarrow\exists x\neg \varphi.
    \]
\end{metacorollary}

\begin{proof}
    Sabemos que $T\vdash \neg\exists x\neg \varphi \leftrightarrow \forall x\neg\neg\varphi$.
    
    Como $T\vdash \neg\neg\varphi\leftrightarrow\varphi$, pela transitividade de $\leftrightarrow$, segue que $T\vdash \forall x\neg\neg\varphi\leftrightarrow\forall x\varphi$.

    Pela transitividade de $\leftrightarrow$, $T\vdash \neg\exists x\neg \varphi \leftrightarrow \forall x\varphi$.
    Novamente pela contrapositiva, $T\vdash \neg\forall x\varphi \leftrightarrow \exists x\neg \varphi$.
\end{proof}


A seguir, precisaremos de uma forma simples de generalização universal.

\begin{metalemma}[Generalização de teoremas]\label{mlem:generalizacaoDeTeoremas}
    Seja $T$ uma teoria de primeira ordem, $\varphi$ uma fórmula da linguagem de $T$ e $x$ uma variável.
    Se $T\vdash \varphi$, então $T\vdash \forall x\varphi$.

    \[
        \frac{T\vdash \varphi}
             {T\vdash \forall x\varphi}.
    \]

    ``Temos que $\varphi$.
    Como $x$ é arbitrário, $\varphi$ vale para todo $x$.''
\end{metalemma}

\begin{proof}
    Escolha uma variável $z$ distinta de $x$.

    Por FO16, $T\vdash z=z$.
    Pela introdução de $\rightarrow$, $T\vdash z=z\rightarrow\varphi$.

    Como $x$ não ocorre livre em $z=z$, podemos aplicar a regra de introdução do quantificador universal e obter
    \[
        T\vdash (z=z)\rightarrow\forall x\varphi.
    \]
    Por \emph{modus ponens} com $T\vdash z=z$, concluímos
    \[
        T\vdash \forall x\varphi.
    \]
\end{proof}

Também usaremos a seguinte renomeação de variável ligada.

\begin{metalemma}[Renomeação de variável universal]\label{mlem:renomeacaoVariavelUniversal}\index{renomeação de variável!universal}
    Seja $T$ uma teoria de primeira ordem.
    Seja $\varphi$ uma fórmula e sejam $x$ e $y$ variáveis.
    Suponha que $y$ não ocorre em $\varphi$.

    Então $T\vdash \forall x\varphi$ se, e somente se, $T\vdash \forall y\,\varphi[x/y]$.
    
    Ademais, para a implicação direta, basta que $x$ seja substituível por $y$ em $\varphi$.

    \[
        \frac{T\vdash \forall x\varphi}
             {T\vdash \forall y\,\varphi[x/y]}
        \qquad
        \frac{T\vdash \forall y\,\varphi[x/y]}
             {T\vdash \forall x\varphi}.
    \]

    ``Temos $\forall x\varphi$.
    Renomeando $x$ para uma variável nova $y$, obtemos $\forall y\,\varphi[x/y]$.''
\end{metalemma}

\begin{proof}
    Suponha primeiro que $T\vdash \forall x\varphi$.

    Como $y$ não ocorre em $\varphi$, a substituição $\varphi[x/y]$ é válida.
    Por instanciação universal,
    \[
        T\vdash \varphi[x/y].
    \]
    Pelo Metalema~\ref{mlem:generalizacaoDeTeoremas}, $T\vdash \forall y\,\varphi[x/y]$.
    Note que só foi necessário que $x$ seja substituível por $y$ em $\varphi$.

    Reciprocamente, suponha que
    \[
        T\vdash \forall y\,\varphi[x/y].
    \]
    Como $y$ não ocorre em $\varphi$, temos que $x$ não ocorre no escopo de nenhum quantificador de $\varphi$ que age em $y$, pois tais quantificadores não ocorrem em $\varphi$.
    Assim, a substituição $\varphi[x/y]$ é válida.

    Pelo já provado, $T\vdash \forall x\,\varphi[x/y][y/x]$.
    Como $\varphi[x/y][y/x]$ é a própria $\varphi$, segue a tese.
\end{proof}

Agora trataremos da prática de introduzir um nome temporário.
Formalmente, isso será feito acrescentando à linguagem um símbolo de constante novo.

\begin{metalemma}[Generalização Universal]\label{mlem:generalizacaoUniversal}
    Seja $T=(\mathcal L,\Sigma)$ uma teoria de primeira ordem.
    Seja $c$ um símbolo de constante que não pertence a $\mathcal L$.
    Seja $\mathcal L(c)$ a linguagem obtida de $\mathcal L$ acrescentando o símbolo $c$, e seja $T(c)=(\mathcal L(c),\Sigma)$.

    Suponha que $\theta$ é uma $\mathcal L)$-fórmula com única variável possivelmente livre $x$, e que $T(c)\vdash \theta$.
    Então $T\vdash \forall x\theta$.

    \[
        \frac{T(c)\vdash \theta}
             {T\vdash \forall x\theta}.
    \]
\end{metalemma}

Antes de demonstrarmos esse resultado, convém discutir sua importância.
A aplicação deste argumento é extremamente corriqueira em matemática.
Ele justifica a prática de introduzir um nome temporário para um elemento arbitrário no formato a seguir.

``Queremos mostrar que $\forall x \theta$.
Fixe um tal $x$, digamos, $c$.
Veremos que $\theta[x/c]$.''

No cotidiano matemático, muitas vezes tal argumento acaba sendo escrito apenas como a seguir.

``Queremos mostrar que $\forall x \theta$.
Fixe $x$.
Veremos que $\theta(x)$.''

Ou seja, na prática, a letra $x$ muitas vezes é inalterada.
Porém, a palavra ``fixe'' dá a entender que a partir daquele ponto do texto, $x$ será temporariamente tratada como um objeto particular, um nome para algum objeto fixado $x$ particular.
A partir desse ponto, $x$ é tratada como uma constante, apesar de mantida a letra $x$.

\begin{proof}
    Seja $\theta_0,\dots,\theta_n$ uma demonstração de $\theta[x/c]$ em $T(c)$, com $\theta_n=\theta[x/c]$.

    Como a demonstração é finita, escolha uma variável $y$ que não ocorre em $\theta$ nem em nenhuma das fórmulas $\theta_0,\dots,\theta_n$.
    Em particular, $y$ não ocorre em nenhum termo que aparece na demonstração.

    Para cada $i\leq n$, denote por $\theta_i^\dagger$ a fórmula obtida de $\theta_i$ substituindo todas as ocorrências do símbolo de constante $c$ pela variável $y$.

    Formalmente, define-se, por recursão na complexidade de cada $L(c)$-termo $t$ no qual $y$ não ocorre, um termo $t^\dagger$ da seguinte forma:
    \begin{itemize}
        \item Se $t$ é uma variável, então $t^\dagger$ é o próprio $t$.
        \item Se $t$ é um símbolo de constante diferente de $c$, então $t^\dagger$ é o próprio $t$.
        \item Se $t$ é o símbolo de constante $c$, então $t^\dagger$ é a variável $y$.
        \item Se $t$ é $f(t_1,\dots,t_k)$, então $t^\dagger$ é $f(t_1^\dagger,\dots,t_k^\dagger)$, onde $f$ é um símbolo de função de aridade $k\geq 1$.
    \end{itemize}

    Então, define-se, para cada $L(c)$-fórmula $\psi$ na qual $y$ não ocorre, a fórmula $\psi^\dagger$ da seguinte forma:
    \begin{itemize}
        \item Se $\psi$ é $t_1=t_2$, então $\psi^\dagger$ é $t_1^\dagger=t_2^\dagger$.
        \item Se $\psi$ é $R(t_1,\dots,t_k)$, então $\psi^\dagger$ é $R(t_1^\dagger,\dots,t_k^\dagger)$, onde $R$ é um símbolo de predicado de aridade $k\geq 1$.
        \item Se $\psi$ é $R$, então $\psi^\dagger$ é o próprio $R$, onde $R$ é um símbolo de predicado de aridade $0$.
        \item Se $\psi$ é $\neg\chi$, então $\psi^\dagger$ é $\neg\chi^\dagger$.
        \item Se $\psi$ é $\chi\square\eta$, então $\psi^\dagger$ é $\chi^\dagger\square\eta^\dagger$, onde $\square$ é um conectivo lógico binário.
        \item Se $\psi$ é $Qz\chi$, então $\psi^\dagger$ é $Qz\chi^\dagger$, onde $Q$ é um quantificador e $z$ é uma variável.
    \end{itemize}

    Precisaremos das seguintes afirmações.

    \begin{claim}\label{claim:termosInalterados}
        Se $t$ é um $L$-termo, então $t^\dagger$ é o próprio $t$.
    \end{claim}
    \begin{proof}
        Procede-se por indução na complexidade de $t$.

        Se $t$ é uma variável, então $t^\dagger$ é o próprio $t$.
        Se $t$ é um símbolo de constante, então tal constante não é $c$, pois $c$ não pertence a $L$.
        Logo, também nesse caso, $t^\dagger$ é o próprio $t$.

        Se $t$ é $f(t_1,\dots,t_k)$, então $t^\dagger=f(t_1^\dagger,\dots,t_k^\dagger)$.
        Pela hipótese de indução, $t_i^\dagger$ é $t_i$ para cada $i\leq k$.
        Assim, $t^\dagger=f(t_1,\dots,t_k)$, isto é, $t^\dagger$ é o próprio $t$.
    \end{proof}

    \begin{claim}\label{claim:formulasInalteradas}
        Se $\psi$ é uma $L$-fórmula tal que $y$ não ocorre em $\psi$, então $\psi^\dagger$ é a própria $\psi$.
    \end{claim}
    \begin{proof}
        Procede-se por indução na complexidade de $\psi$.

        Se $\psi$ é $t_1=t_2$, $R(t_1,\dots,t_k)$ ou $R$, então $\psi^\dagger$ é a própria $\psi$, pois os termos envolvidos são $L$-termos e, portanto, são inalterados pela Afirmação~\ref{claim:termosInalterados}.

        Se $\psi$ é $\neg\chi$, então $\psi^\dagger$ é $\neg\chi^\dagger$, que, por hipótese de indução, é $\neg\chi$, ou seja, a própria $\psi$.

        Se $\psi$ é $\chi\square\eta$, então $\psi^\dagger$ é $\chi^\dagger\square\eta^\dagger$, que, por hipótese de indução, é $\chi\square\eta$, ou seja, a própria $\psi$.

        Se $\psi$ é $Qz\chi$, então $\psi^\dagger$ é $Qz\chi^\dagger$, que, por hipótese de indução, é $Qz\chi$, ou seja, a própria $\psi$.
    \end{proof}

    \begin{claim}\label{claim:termosComY}
        Seja $t$ um $L(c)$-termo tal que $y$ não ocorre em $t$.
        Então toda variável de $t^\dagger$ é ou uma variável de $t$ ou a variável $y$.
    \end{claim}
    \begin{proof}
        Procede-se por indução na complexidade de $t$.

        Se $t$ é uma variável, digamos, $z$, então $t^\dagger$ é o próprio $z$.
        Assim, a única variável de $t^\dagger$ é uma variável de $t$.

        Se $t$ é um símbolo de constante diferente de $c$, então $t^\dagger$ é o próprio $t$, que não possui variáveis.

        Se $t$ é o símbolo de constante $c$, então $t^\dagger$ é a variável $y$.

        Se $t$ é $f(t_1,\dots,t_k)$, então $t^\dagger=f(t_1^\dagger,\dots,t_k^\dagger)$.
        As variáveis de $t^\dagger$ são as variáveis dos termos $t_1^\dagger,\dots,t_k^\dagger$.
        Pela hipótese de indução, cada uma delas é uma variável de algum $t_i$ ou é a variável $y$.
        Logo, cada variável de $t^\dagger$ é uma variável de $t$ ou é a variável $y$.
    \end{proof}

    \begin{claim}\label{claim:variaveisLivresComY}
        Seja $\varphi$ uma $L(c)$-fórmula tal que $y$ não ocorre em $\varphi$.
        Toda variável livre de $\varphi^\dagger$ é ou uma variável livre de $\varphi$ ou a variável $y$.
    \end{claim}
    \begin{proof}
        Procede-se por indução na complexidade de $\varphi$.

        Se $\varphi$ é $t_1=t_2$, $R(t_1,\dots,t_k)$ ou $R$, então o resultado segue diretamente da Afirmação~\ref{claim:termosComY}.

        Se $\varphi$ é $\neg\chi$, então $\varphi^\dagger$ é $\neg\chi^\dagger$.
        As variáveis livres de $\varphi^\dagger$ são as variáveis livres de $\chi^\dagger$.
        Pela hipótese de indução, estas são variáveis livres de $\chi$ ou a variável $y$.
        Portanto, são variáveis livres de $\varphi$ ou a variável $y$.

        Se $\varphi$ é $\chi\square\eta$, então $\varphi^\dagger$ é $\chi^\dagger\square\eta^\dagger$.
        As variáveis livres de $\varphi^\dagger$ são as variáveis livres de $\chi^\dagger$ ou de $\eta^\dagger$.
        Pela hipótese de indução, estas são variáveis livres de $\chi$ ou de $\eta$, ou a variável $y$.
        Portanto, são variáveis livres de $\varphi$ ou a variável $y$.

        Se $\varphi$ é $Qz\chi$, então $\varphi^\dagger$ é $Qz\chi^\dagger$.
        As variáveis livres de $\varphi^\dagger$ são as variáveis livres de $\chi^\dagger$ diferentes de $z$.
        Pela hipótese de indução, estas são variáveis livres de $\chi$ ou a variável $y$.
        Como $y$ não ocorre em $\varphi$, em particular $z$ não é $y$.
        Logo, as variáveis livres de $\varphi^\dagger$ são variáveis livres de $\varphi$ ou a variável $y$.
    \end{proof}

    \begin{claim}\label{claim:substituicaoValidaDagger}
        Seja $\varphi$ uma $L(c)$-fórmula tal que $y$ não ocorre em $\varphi$, seja $t$ um $L(c)$-termo tal que $y$ não ocorre em $t$, e seja $z$ uma variável.
        Se $z$ é substituível por $t$ em $\varphi$, então $z$ é substituível por $t^\dagger$ em $\varphi^\dagger$.
    \end{claim}
    \begin{proof}
        Considere uma ocorrência livre de $z$ em $\varphi^\dagger$.
        Como a operação $(\cdot)^\dagger$ apenas substitui o símbolo de constante $c$ pela variável $y$, essa ocorrência de $z$ já correspondia a uma ocorrência livre de $z$ em $\varphi$.

        Suponha que essa ocorrência esteja no escopo de um quantificador que age em uma variável $w$ que ocorre em $t^\dagger$.
        Como $y$ não ocorre em $\varphi$, a fórmula $\varphi^\dagger$ não possui quantificadores agindo em $y$.
        Portanto, $w$ não é $y$.

        Pela Afirmação~\ref{claim:termosComY}, como $w$ ocorre em $t^\dagger$ e $w\neq y$, então $w$ ocorre em $t$.
        Assim, em $\varphi$, a ocorrência livre correspondente de $z$ estava no escopo de um quantificador que age em uma variável que ocorre em $t$, contrariando a hipótese de que $z$ é substituível por $t$ em $\varphi$.

        Logo, nenhuma ocorrência livre de $z$ em $\varphi^\dagger$ está no escopo de um quantificador que age em uma variável que ocorre em $t^\dagger$.
        Portanto, $z$ é substituível por $t^\dagger$ em $\varphi^\dagger$.
    \end{proof}

    \begin{claim}\label{claim:daggerComutaSubstituicao}
        Seja $\varphi$ uma $L(c)$-fórmula tal que $y$ não ocorre em $\varphi$, seja $z$ uma variável e seja $t$ um $L(c)$-termo tal que $y$ não ocorre em $t$.
        Então
        \[
            (\varphi[z/t])^\dagger=\varphi^\dagger[z/t^\dagger].
        \]
    \end{claim}
    \begin{proof}
        Primeiro, provamos por indução na complexidade dos termos $s$ que
        \[
            (s[z/t])^\dagger=s^\dagger[z/t^\dagger].
        \]

        Se $s$ é a variável $z$, então $s[z/t]$ é $t$, logo $(s[z/t])^\dagger=t^\dagger$.
        Por outro lado, $s^\dagger$ é $z$, e $z[z/t^\dagger]$ é $t^\dagger$.

        Se $s$ é uma variável diferente de $z$, então $s[z/t]$ é $s$, e ambos os lados são iguais a $s$.

        Se $s$ é um símbolo de constante, então $s[z/t]$ é $s$, e ambos os lados são iguais a $s^\dagger$.

        Se $s$ é $f(s_1,\dots,s_k)$, então, pela hipótese de indução aplicada aos $s_i$,
        \[
            (s[z/t])^\dagger
            =
            f((s_1[z/t])^\dagger,\dots,(s_k[z/t])^\dagger)
            =
            f(s_1^\dagger[z/t^\dagger],\dots,s_k^\dagger[z/t^\dagger])
            =
            s^\dagger[z/t^\dagger].
        \]

        Agora procedemos por indução na complexidade de $\varphi$.

        Se $\varphi$ é $s_1=s_2$, $R(s_1,\dots,s_k)$ ou $R$, a conclusão segue do resultado já provado para termos.

        Se $\varphi$ é $\neg\chi$, então
        \[
            (\varphi[z/t])^\dagger
            =
            \neg(\chi[z/t])^\dagger
            =
            \neg(\chi^\dagger[z/t^\dagger])
            =
            \varphi^\dagger[z/t^\dagger],
        \]
        pela hipótese de indução.

        Se $\varphi$ é $\chi\square\eta$, então a conclusão segue imediatamente da hipótese de indução aplicada a $\chi$ e a $\eta$.

        Se $\varphi$ é $Qw\chi$, com $w$ diferente de $z$, então
        \[
            (\varphi[z/t])^\dagger
            =
            Qw(\chi[z/t])^\dagger
            =
            Qw(\chi^\dagger[z/t^\dagger])
            =
            \varphi^\dagger[z/t^\dagger].
        \]

        Se $\varphi$ é $Qz\chi$, então $\varphi[z/t]$ é $\varphi$.
        Logo, $(\varphi[z/t])^\dagger=\varphi^\dagger$.
        Por outro lado, como a substituição por $z$ não atravessa um quantificador que age em $z$, temos $\varphi^\dagger[z/t^\dagger]=\varphi^\dagger$.
        Portanto, $(\varphi[z/t])^\dagger=\varphi^\dagger[z/t^\dagger]$.
    \end{proof}

    \begin{claim}\label{claim:substituirConstantePorVariavel}
        Seja $\varphi$ uma $L$-fórmula tal que $y$ não ocorre em $\varphi$, e seja $z$ uma variável.
        Então $(\varphi[z/c])^\dagger=\varphi[z/y]$.
    \end{claim}
    \begin{proof}
        Pela Afirmação~\ref{claim:daggerComutaSubstituicao}, aplicada ao termo $c$, $(\varphi[z/c])^\dagger=\varphi^\dagger[z/c^\dagger]$.

        Pela Afirmação~\ref{claim:formulasInalteradas}, como $\varphi$ é uma $L$-fórmula e $y$ não ocorre em $\varphi$, temos $\varphi^\dagger=\varphi$.
        Além disso, $c^\dagger$ é $y$.
        Logo, $(\varphi[z/c])^\dagger=\varphi[z/y]$.
    \end{proof}

    Provaremos, por indução em $i\leq n$, que $T\vdash \theta_i^\dagger$.
    Suponha que a tese já tenha sido provada para todo $j<i$.
    Veremos que $T\vdash \theta_i^\dagger$.

    Se $\theta_i$ é um axioma de $T(c)$, então $\theta_i\in\Sigma$.
    Como os axiomas de $\Sigma$ são $L$-sentenças, $\theta_i$ é uma $L$-fórmula.
    Como $y$ não ocorre em $\theta_i$, pela Afirmação~\ref{claim:formulasInalteradas}, $\theta_i^\dagger$ é a própria $\theta_i$.
    Assim, $T\vdash \theta_i^\dagger$.

    Se $\theta_i$ é um axioma lógico de $L(c)$, então $\theta_i^\dagger$ é um axioma lógico de $L$.
    De fato, como a operação $(\cdot)^\dagger$ comuta com os conectivos lógicos, cada instância de um dos esquemas FO1--FO13 é levada a uma instância do mesmo esquema.
    O axioma FO16 é preservado imediatamente.

    Quanto a FO17 e FO18, note que o novo símbolo $c$ é apenas um símbolo de constante.
    Portanto, ele não é um símbolo funcional de aridade positiva nem um símbolo de predicado, e não introduz novos casos nesses esquemas.
    Assim, instâncias de FO17 e FO18 em $L(c)$ são levadas a instâncias dos mesmos esquemas em $L$.

    Resta verificar FO14 e FO15.
    Se $\theta_i$ é uma instância de FO14, então $\theta_i$ é da forma $\forall z\varphi\rightarrow\varphi[z/t]$ em que $z$ é substituível por $t$ em $\varphi$.
    Como $y$ não ocorre em $\theta_i$, temos que $y$ não ocorre em $\varphi$.
    Podemos assumir que $y$ não ocorre em $t$, pois, caso contrário, como a substituição $\varphi[z/t]$ é válida, teríamos que $z$ não ocorre livre em $\varphi$ e, portanto, tal axioma é o mesmo que $\varphi\rightarrow\varphi[z/z]$.

    Pelas Afirmações~\ref{claim:substituicaoValidaDagger} e~\ref{claim:daggerComutaSubstituicao}, temos que $z$ é substituível por $t^\dagger$ em $\varphi^\dagger$ e $(\varphi[z/t])^\dagger=\varphi^\dagger[z/t^\dagger]$.
    Portanto, $\theta_i^\dagger$ é $\forall z\varphi^\dagger\rightarrow\varphi^\dagger[z/t^\dagger]$, que é uma instância de FO14 em $L$.

    O caso de FO15 é tratado de forma análoga.

    Suponha agora que $\theta_i$ foi obtida por \emph{modus ponens}.
    Então existem $j,k<i$ tais que $\theta_j$ é $\theta_k\rightarrow\theta_i$.
    Pela hipótese de indução,
    \[
        T\vdash \theta_j^\dagger
        \qquad\text{e}\qquad
        T\vdash \theta_k^\dagger.
    \]
    Como $\theta_j^\dagger$ é $\theta_k^\dagger\rightarrow\theta_i^\dagger$, segue por \emph{modus ponens} que $T\vdash \theta_i^\dagger$.

    Suponha que $\theta_i$ foi obtida pela regra de introdução do quantificador universal.
    Então existem fórmulas $\chi$ e $\eta$, uma variável $z$ e $j<i$ tais que $\theta_j$ é $\chi\rightarrow\eta$, $\theta_i$ é $\chi\rightarrow\forall z\eta$, e $z$ não ocorre livre em $\chi$.

    Pela hipótese de indução, $T\vdash \chi^\dagger\rightarrow\eta^\dagger$.
    Como $y$ não ocorre em $\theta_i$, temos $z\neq y$.
    Pela Afirmação~\ref{claim:variaveisLivresComY}, toda variável livre de $\chi^\dagger$ é uma variável livre de $\chi$ ou é $y$.
    Como $z$ não ocorre livre em $\chi$ e $z\neq y$, segue que $z$ não ocorre livre em $\chi^\dagger$.

    Pela regra de introdução do quantificador universal,
    \[
        T\vdash \chi^\dagger\rightarrow\forall z\eta^\dagger.
    \]
    Essa fórmula é exatamente $\theta_i^\dagger$.

    Por fim, suponha que $\theta_i$ foi obtida pela regra de introdução do quantificador existencial.
    Então existem fórmulas $\chi$ e $\eta$, uma variável $z$ e $j<i$ tais que $\theta_j$ é $\chi\rightarrow\eta$, $\theta_i$ é $\exists z\chi\rightarrow\eta$, e $z$ não ocorre livre em $\eta$.

    Pela hipótese de indução, $T\vdash \chi^\dagger\rightarrow\eta^\dagger$.
    Como $y$ não ocorre em $\theta_i$, temos $z\neq y$.
    Pela Afirmação~\ref{claim:variaveisLivresComY}, toda variável livre de $\eta^\dagger$ é uma variável livre de $\eta$ ou é $y$.
    Como $z$ não ocorre livre em $\eta$ e $z\neq y$, segue que $z$ não ocorre livre em $\eta^\dagger$.

    Pela regra de introdução do quantificador existencial,
    \[
        T\vdash \exists z\chi^\dagger\rightarrow\eta^\dagger.
    \]
    Essa fórmula é exatamente $\theta_i^\dagger$.

    Isso encerra a indução.
    Em particular, para $i=n$, obtemos $T\vdash (\theta[x/c])^\dagger$.
    Pela Afirmação~\ref{claim:substituirConstantePorVariavel}, segue que
    \[
        T\vdash \theta[x/y].
    \]

    Pelo Metalema~\ref{mlem:generalizacaoDeTeoremas},
    \[
        T\vdash \forall y\,\theta[x/y].
    \]
    Como $y$ não ocorre em $\theta$, pelo Metalema~\ref{mlem:renomeacaoVariavelUniversal}, concluímos que
    \[
        T\vdash \forall x\,\theta.
    \]
\end{proof}

Temos uma equivalênca

\begin{metatheorem}[Instanciação existencial]\label{mthm:instanciacaoExistencial}
    Seja $T=(\mathcal L,\Sigma)$ uma teoria de primeira ordem, $x$ uma variável, e $\varphi$ uma fórmula da linguagem de $T$ tal que $x$ é a única variável possivelmente livre.
    Seja $\psi$ uma fórmula de $\mathcal L$ sem $x$ livre.


    Seja $c$ um símbolo de constante que não pertence a $\mathcal L$, e seja $\mathcal L(c)$ a linguagem obtida de $\mathcal L$ acrescentando o símbolo $c$.
    Seja $T(c)=(\mathcal L(c),\Sigma)$.
    
    Se $T(c)\vdash \varphi[x/c]\rightarrow\psi$, então $T\vdash \exists x\varphi\rightarrow\psi$.

    \[
       \frac{T(c)\vdash \varphi[x/c]\rightarrow\psi}
             {T\vdash \exists x\varphi\rightarrow\psi}.
    \]
\end{metatheorem}

Novamente, antes de demonstrar esse resultado, convém discutir sua importância.
É corriqueiro em matemática a necessidade de se demonstrar que se existe um objeto $x$ que satisfaz determinada propriedade $\varphi$, então alguma outra propriedade $\psi$ é satisfeita.
Nessas hipóteses, o argumento a ser seguido é o seguinte.

``Queremos mostrar que se existe um $x$ tal que $\varphi$, então $\psi$.
Fixe um tal $x$, digamos, $c$ (...) assim, $\psi$.
Logo, se existe um $x$ tal que $\varphi$, então $\psi$.''

Note que, na prática, o nome temporário $c$ é muitas vezes escrito como $x$, e a letra $x$ é mantida.
Porém, como já discutido na discussão sobre a generalização universal, a palavra ``fixe'' dá a entender que a partir daquele ponto do texto, $x$ será temporariamente tratada como constante.

\begin{proof}
    Suponha que $T(c)\vdash \varphi[x/c]\rightarrow\psi$.
    Como $\psi$ não possui $x$ livre, então $(\varphi\rightarrow \psi)[x/c]$ é $\varphi[x/c]\rightarrow\psi$.
    Assim, $T(c)\vdash (\varphi\rightarrow \psi)[x/c]$.
    Pelo Metalema~\ref{mlem:generalizacaoUniversal}, $T\vdash \forall x(\varphi\rightarrow \psi)$.
    
    Como a substituição de $x$ por $x$ é válida, segue da instanciação universal que $T\vdash \varphi\rightarrow \psi$.
    Pela regra de introdução do quantificador existencial, $T\vdash \exists x\varphi\rightarrow \psi$.
\end{proof}
